\section{Related Literature}\label{sec:literature}

\paragraph{Simulation optimization and noisy constraints.}
Simulation optimization studies expensive stochastic systems for which
objectives or constraints are available only through simulation
\citep{Fu2002,AmaranEtAl2016}.  Stochastic kriging explicitly separates
response uncertainty from simulation noise \citep{AnkenmanNelsonStaum2010},
while constrained Bayesian optimization models noisy objective and feasibility
observations jointly \citep{LethamEtAl2019}.  Our question occurs one stage
earlier: before fitting a target surrogate, which few ordered profiles should
receive the first simulations?

\paragraph{High-dimensional and functional optimization.}
Trust-region methods and sparse axis-aligned priors can make Bayesian
optimization more robust in large ambient spaces
\citep{ErikssonEtAl2019,ErikssonPoloczek2021,ErikssonJankowiak2021}.  They do
not by themselves encode that coordinates discretize one ordered function.
Functional Bayesian optimization instead searches over function-valued
decisions \citep{VienEtAl2018}.  Functional-data analysis likewise represents
curves through common bases and treats discretization as observation of a
function rather than unrelated coordinates \citep{RamsaySilverman2005}.  Our
profile coordinate is less flexible but
fully finite: it maps irregular or refined grids into a common cosine summary,
then selects from a declared structural library.  Nominal grid dimension and
effective profile complexity are therefore reported separately.

\paragraph{Space-filling initial designs.}
Maximin and minimax distance criteria are classical model-free devices for
covering an experimental region \citep{JohnsonMooreYlvisaker1990}, and
space-filling designs are standard for deterministic computer experiments
\citep{MorrisMitchell1995}.  Gonzalez farthest-first gives a factor-two
finite $k$-center guarantee \citep{Gonzalez1985}.  We use that algorithm on a
finite profile library, but source ranks augment its selection metric.  The
standardized DCT maximin control uses the same structural coordinates without
source ranks.  It separates source information from the coordinate scaling
used by the original unstandardized DCT control.

\paragraph{Learned initializations and portfolios.}
Learning starting configurations from historical tasks is established work.
\citet{WistubaEtAl2015} optimize a meta-loss to learn an initialization
sequence, while \citet{FeurerEtAl2022} greedily construct a finite portfolio
of complementary configurations and evaluate it before continuing Bayesian
optimization.  Our contribution concerns ordered policies under noisy chance
constraints: a common structural coordinate across grids, joint safety and
objective source ranks, and independent deployment verification with explicit
source costs.  To compare selection rules within this setting, we include a
source-greedy portfolio using the same archive and ten target evaluations.

\paragraph{Transfer and historical-task priors.}
Few-shot deep kernels, hierarchical Gaussian-process priors, source ensembles,
and meta-learned acquisitions transfer different objects across tasks
\citep{WistubaGrabocka2021,RothfussEtAl2021,RothfussEtAl2021FPACOH,RothfussEtAl2023,
WangEtAl2024HyperBO,TighineanuEtAl2022,VolppEtAl2020,PanEtAl2024}.
Related work also learns a search space from earlier hyperparameter tasks
\citep{PerroneEtAl2019} or weights source surrogates by ranking performance
\citep{FeurerEtAl2018}.  Warm-start Bayesian optimization similarly treats
earlier evaluations as reusable information \citep{PoloczekWangFrazier2016}.
Our method transfers a finite initial design: the source archive ranks a fixed
public library, while the target surrogate and optimizer remain separate.  This makes
it possible to compare source scoring directly with the same structural
library without outcomes and to count the cost of history explicitly.  We
also compare eight end-to-end transfer pipelines under their own common
source, target-budget, and verification protocol.

\paragraph{Selection and independent verification.}
Chance-constrained ranking-and-selection procedures allocate simulation to
identify a feasible best system while controlling selection error
\citep{HongLuoNelson2015}.  Post-optimization ranking and selection can also
``clean up'' candidate solutions generated by a heuristic search
\citep{BoeselNelsonKim2003}.  We adopt this separation literally.  Search
produces an ordered shortlist; fresh Bernoulli feasibility replications test
its members; and Bonferroni spending composes candidate-wise exact tests.  The
all-success rule is conservative, so we report power and verification cost in
addition to false certificates.  This differs from treating a surrogate's
posterior feasibility probability as a deployment certificate.

\paragraph{Positioning.}
Without restrictions on the objective and feasible set, no finite
target-label-free design can dominate all others, a point consistent with
no-free-lunch results \citep{WolpertMacready1997}.  Our claim is therefore
about experimental design for related ordered-profile tasks, not a
distribution-free optimizer.  Source-free structured designs, adverse
regimes, an external negative control, and explicit cost accounting show where
that claim does and does not hold.
